\documentclass[sigconf,nonacm]{acmart}

%% ACM 메타데이터 간소화
\settopmatter{printacmref=false}
\renewcommand\footnotetextcopyrightpermission[1]{}
\pagestyle{plain}

%% 한글 지원 (XeLaTeX + kotex 권장)
\usepackage{kotex}

%% 기타 패키지
\usepackage{booktabs}
\usepackage{graphicx}
\usepackage{hyperref}
\usepackage{listings}
\usepackage{xcolor}

%% 긴 식별자/경로가 좁은 2-column에서 margin을 벗어나는 문제 방지.
%% url 패키지(hyperref 내부)로 타이프라이터 스타일의 breakable 명령 \code 정의.
\makeatletter
\g@addto@macro\UrlBreaks{\do\_\do\.\do\/\do\-\do\(\do\)\do\,\do\=\do\"}
\makeatother
\DeclareUrlCommand\code{\urlstyle{tt}}

%% 전반적으로 줄바꿈을 느슨하게
\sloppy

\begin{document}

\title{Jamo-based Korean Tokenization: \\
A BPE Approach for Sub-character Level Subword Units}

\author{Minjae Kyung}
\affiliation{%
  \institution{Digipen Institute of Technology}
  \country{USA}
}
\email{minjae.kyung@digipen.com}

\author{Ju-ve Chankasemporn}
\affiliation{%
  \institution{Digipen Institute of Technology}
  \country{USA}
}
\email{juve.chankasemporn@digipen.edu}


\begin{abstract}
현대 자연어 처리에서 토크나이저는 모델 학습과 추론 성능에 직접적인 영향을 미치는 핵심 구성 요소이다.
그러나 대부분의 다국어 토크나이저는 문자(character)를 최소 단위로 가정하고 Byte Pair Encoding (BPE)을 적용하기 때문에,
Hangul과 같이 문자 자체가 자음과 모음의 결합으로 구성되는 표기 체계에서는
형태소 경계가 음절 내부에 존재하여도 이를 포착하지 못하는 구조적 한계가 있다.
본 연구에서는 Hangul 음절을 초성(initial), 중성(medial), 종성(final)의 Jamo 단위로 분해한 뒤
BPE를 적용하는 Jamo 기반 토크나이저를 제안한다. 자체 구현한 Python BPE와
Hugging Face \texttt{tokenizers}~\cite{huggingface_tokenizers} 기반 BPE의 두 버전을 32k 및 16k vocab
크기로 학습하였으며, 동일 구조의 Transformer~\cite{vaswani2017attention} 기반 한영 번역기를 두 토크나이저로
각각 학습시켜 토크나이저가 학습에 미치는 영향을 분석하였다. CC-100 한국어 코퍼스~\cite{conneau2020cc100}에 대한
토큰화 통계 실험 결과, Jamo 기반 토크나이저는 동일한 vocab 크기에서 문자 기반 대비 평균 토큰 수가 더 적었으며,
특히 도메인 한정 소규모 코퍼스(소설 데이터)로 학습한 경우와 작은 vocab 크기에서 그 격차가 확대되었다. 또한 ``노랫말''과 같이 형태소 경계가
음절 내부에 위치하는 사례에서 Jamo 기반 분절이 의미 단위와 일치하는 결과를 보였다.
\end{abstract}

\keywords{Korean Tokenization, Jamo, Byte Pair Encoding, Subword Tokenization, Neural Machine Translation, Hangul}

\maketitle

\section{Introduction}
\label{sec:intro}

자연어 처리(NLP) 모델의 성능은 입력 텍스트를 얼마나 효과적으로 토큰 단위로 분절하는가에 크게 의존한다.
현대의 대규모 언어 모델과 기계 번역 시스템은 대부분 Byte Pair Encoding (BPE)~\cite{sennrich2016bpe},
WordPiece~\cite{wu2016google}, SentencePiece~\cite{kudo2018sentencepiece}와 같은 서브워드 기반
토크나이저를 사용한다. 이러한 토크나이저는 공통적으로 \emph{문자(character)} 또는 \emph{바이트(byte)}를
최소 단위로 가정하고, 자주 함께 등장하는 인접 기호 쌍을 반복적으로 병합하여 서브워드 어휘를 구성한다.

그러나 이러한 가정은 모든 언어에 보편적으로 적용되지 않는다. 한국어의 표기 체계인
Hangul은 자음(초성, 종성)과 모음(중성)이 결합하여 하나의 음절 블록을 이루는 특수한 구조를 갖는다.
예를 들어 ``먹다''의 ``먹''은 시각적으로는 단일 문자이지만 실제로는 $\langle$ㅁ, ㅓ, ㄱ$\rangle$의
세 Jamo로 구성되며, 형태소나 음운 현상의 경계가 음절 \emph{내부}에 존재하는 경우가 드물지 않다.
문자 단위에서만 토큰을 병합하면 Hangul 음절을 최소 원자로 취급하게 되어, 이처럼 음절 내부에 걸친
의미 경계를 표현할 수 없다.

이러한 한계는 실무적으로도 이미 관측되고 있다. GPT 계열의 tiktoken과 같은 다국어 토크나이저에서는
한국어 한 문장이 같은 뜻의 영어 문장보다 수 배 많은 토큰을 차지하는 경우가 흔하며, 이는 API
호출 비용, 문맥 길이 활용률, 추론 지연(latency) 등 다운스트림 효용에 직접적인 손실을 야기한다.
Hangul의 미세 구조(fine-grained structure)를 활용하면 이 손실을 부분적으로나마 보전할 수 있으리라
기대된다.

본 연구는 이러한 문제를 해결하기 위해 다음의 기여를 제시한다:
\begin{itemize}
    \item Hangul 음절을 결정론적 공식을 통해 Jamo 단위로 분해한 뒤 BPE를 적용하는
    Jamo 기반 토크나이저를 구현하였다. 디코딩 시에는 Jamo 시퀀스를 다시 음절로 재조합한다.
    \item 교육적 목적의 \textbf{자체 Python 구현}과 대규모 학습용 \textbf{Hugging Face \texttt{tokenizers}
    기반 구현}의 두 가지 버전을 제공하였으며, vocab 크기 32k와 16k의 두 설정으로 각각 학습하였다.
    \item 동일 구조의 Transformer 기반 한영 번역기를 Jamo 기반 / 문자 기반 토크나이저로 각각 학습시켜,
    토크나이저 선택이 실제 학습 곡선(loss curve)에 미치는 영향을 직접 비교하였다.
    \item CC-100 한국어 코퍼스의 한 청크(약 5,000개 문장)에 대해 평균/중앙값/표준편차/최솟값/최댓값 등
    토큰 길이 분포 통계를 측정하여, tiktoken \texttt{o200k\_base}~\cite{tiktoken}을 포함한 기존
    토크나이저들과 정량 비교하였다.
    \item 사용자가 입력 문장의 분절과 번역 결과를 실시간으로 체험할 수 있는 Streamlit 기반 데모를 구축하였다.
\end{itemize}

본 논문의 구성은 다음과 같다. 2장에서는 Hangul의 구조와 BPE 알고리즘에 대한 배경지식을 정리하고,
3장에서는 관련 연구를 개관한다. 4장에서 연구 동기를 상세히 서술한 뒤, 5장에서 제안 방법과 구현
세부사항을, 6장에서 실험 설정을 기술한다. 7장에서 정량/정성 결과를 제시하고, 8장 논의,
9장 한계, 10장 향후 연구에 이어 11장에서 결론을 맺는다.

\section{Background}
\label{sec:background}

\subsection{Hangul and Jamo}

Hangul은 초성(initial consonant), 중성(medial vowel), 선택적 종성(final consonant)의 세 요소가 결합하여
하나의 음절 블록(syllable block)을 형성하는 문자 체계이다. 유니코드에서 현대 Hangul 음절은
\texttt{U+AC00}(가)부터 \texttt{U+D7A3}(힣)까지 총 $19 \times 21 \times 28 = 11{,}172$자로
정의되며, 각 음절 $S$는 다음의 공식을 통해 구성 Jamo로 분해될 수 있다:
\[
S = ((C_{\text{init}} \times 21) + C_{\text{med}}) \times 28 + C_{\text{fin}} + \text{0xAC00}
\]
여기서 $C_{\text{init}} \in [0, 19)$, $C_{\text{med}} \in [0, 21)$, $C_{\text{fin}} \in [0, 28)$은
각각 19개의 초성, 21개의 중성, 28개의 종성(``없음''을 포함) 인덱스이다.
이러한 체계적 규칙 덕분에 한국어는 외부 형태소 분석기 없이도 순수하게 문자열 연산만으로
Jamo 수준까지 결정론적 분해 및 재조합이 가능하다.

본 연구의 구현에서는 Kss~\cite{kss} 라이브러리가 제공하는 \texttt{h2hcj} (Hangul-to-Hangul
Compatibility Jamo) 및 \texttt{hcj2h} 변환을 이용한다. \texttt{h2hcj}는 음절을 호환 자모
(compatibility Jamo, \texttt{U+3131--U+318E}) 시퀀스로 분해하며, \texttt{hcj2h}는 그 역변환을
수행한다. 호환 자모를 사용하는 이유는 초성/종성이 동일한 글리프를 공유하도록 하여
어휘 크기의 낭비를 줄이기 위함이다. 예를 들어 음절 ``강''은 표준 Hangul Jamo 블록
(\texttt{U+1100--U+11FF})에서는 초성 \texttt{U+1100}, 중성 \texttt{U+1161}, 종성 \texttt{U+11BC}의
세 서로 다른 코드포인트로 분해되어 동일한 ``ㄱ''/``ㅇ'' 계열 자모라도 위치에 따라 별개의 심볼로
취급되지만, 호환 자모(\texttt{U+3131--U+318E})에서는 ㄱ (\texttt{U+3131}), ㅏ (\texttt{U+314F}),
ㅇ (\texttt{U+3147})으로 표현되어 위치에 무관하게 동일한 기저 심볼로 통합된다. BPE의 관점에서 이는 merge 통계가 위치 정보에 분산되지 않고 통합된
자모 단위에서 축적됨을 뜻하며, 결과적으로 merge 품질과 vocab 활용 효율을 높인다.

Hangul은 1443년 세종대왕에 의해 창제된 이래 음소(phoneme) 단위를 시각적 블록으로 응집한
음소문자(featural alphabet)로 설계되었다. 이러한 설계는 전산 처리 관점에서
매우 우호적인 특성을 지니는데, (i) 모든 현대 음절이 완전한 조합형 블록으로 정의되어 분해/재조합이
결정론적이고 무손실이며, (ii) 초성/중성/종성의 집합이 작아 vocab 예산 소비가 적으며,
(iii) 고어나 방언 처리를 제외하면 예외 케이스가 거의 없다는 점이다. 본 연구는 이러한
구조적 이점을 적극 활용한다.

\subsection{Byte Pair Encoding}

BPE~\cite{sennrich2016bpe}는 원래 Gage(1994)가 제안한 데이터 압축 기법이었으나,
Sennrich 등에 의해 서브워드 토큰화 분야에 도입되면서 현대 NLP의 사실상 표준 기법 중 하나로
자리잡았다. 알고리즘은 다음과 같이 요약된다.
\begin{enumerate}
    \item 각 단어를 최소 단위(일반적으로 문자) 시퀀스로 분해하고, 단어 끝을 표시하기 위한 특수
    심볼 (예: \texttt{</w>})을 부착한다.
    \item 전 코퍼스에서 가장 빈번한 인접 심볼 쌍을 선택하고, 해당 쌍을 하나의 새로운 심볼로 병합한다.
    \item 원하는 vocab 크기에 도달하거나 최소 빈도 임계값을 만족하는 쌍이 없을 때까지 이를 반복한다.
\end{enumerate}
BPE는 학습 데이터로부터 통계적으로 빈출되는 서브워드를 자연스럽게 포착하기 때문에,
희귀어(rare word) 처리와 사전 크기 제어 측면에서 큰 장점을 지닌다. 기존 연구들은 대부분
문자 혹은 바이트 수준에서 출발하며, Jamo 수준에서 BPE를 적용하는 접근은 비교적 덜 탐구되어 왔다.

BPE의 변형으로는 (i) 모든 유효한 분절 후보에 대해 unigram 언어모델을 이용해 log-likelihood를
최대화하는 유닛을 유지하는 Unigram LM~\cite{kudo2018sp}, (ii) 여러 분절 후보를 확률적으로
sampling하여 정규화 효과를 얻는 BPE-dropout~\cite{provilkov2020bpedropout},
(iii) 바이트 수준에서 출발해 Unicode 정규화에 대한 강건성을 확보하는 byte-level BPE
(GPT-2 이후의 표준)~\cite{radford2019gpt2} 등이 있다. 본 연구는 이러한 변형과 직교하는 축인
\emph{기저 단위}의 선택에 초점을 맞춘다. 즉, 문자/바이트가 아닌 ``자모''라는 더 미세한 기저 단위
위에서 BPE를 수행하는 것의 효과를 측정한다.

\subsection{Pre-tokenization}

BPE는 일반적으로 공백, 문장부호, CJK 범위 등을 기준으로 문자열을 사전에 분리하는 pre-tokenization
단계를 거친다. 본 연구의 HF 구현은 \code{pre_tokenizers.Whitespace()}를 채택하여,
공백과 일반적인 구두점을 경계로 삼는 전통적인 방식을 사용한다. Jamo 모드에서는 Jamo 분해가
pre-tokenization \emph{이전}에 적용되므로, BPE 학습은 공백으로 분리된 Jamo 시퀀스에 대해
진행된다. 이 단계적 구성은 한 음절이 여러 단어에 걸쳐 병합되는 것을 구조적으로 방지한다.

\section{Related Work}
\label{sec:related}

\subsection{Subword Tokenization}

신경 기계 번역과 대규모 언어 모델의 부상 이후 서브워드 토큰화는 사실상 표준 전처리 기법으로
자리 잡았다. Sennrich \emph{et~al.}~\cite{sennrich2016bpe}은 BPE를 최초로 NMT에 적용하여
희귀 어휘와 열린 어휘(open-vocabulary) 문제를 해결할 수 있음을 보였다.
Wu \emph{et~al.}~\cite{wu2016google}은 WordPiece 토크나이저를 Google NMT 시스템에서 제안하여
가능도(likelihood)에 기반한 서브워드 선택을 도입하였고, Kudo \emph{et~al.}~\cite{kudo2018sentencepiece}은
언어 독립적이고 학습 가능한 SentencePiece 및 unigram LM 기반 토크나이저를 공개하였다.
이러한 접근들은 공통적으로 문자 혹은 바이트를 최소 단위로 가정한다.

\subsection{Korean Tokenization}

한국어 토큰화는 주로 두 방향으로 연구되어 왔다.
첫째는 \textbf{형태소 기반 접근}으로, MeCab-ko~\cite{mecabko}, Komoran, Mecab, Okt 등의 형태소 분석기를
이용해 어절을 형태소로 분해한 뒤 서브워드 토큰화를 수행한다. 이는 의미 경계에 정렬된 분절을
제공하지만, 분석기 의존성과 OOV(out-of-vocabulary) 문제를 안는다.
둘째는 \textbf{자모 기반 접근}으로, 일부 연구에서는 Hangul 음절을 초성/중성/종성으로 분해한 뒤
RNN이나 Transformer에 직접 투입하여 문자 수준 모델링의 일반화 성능을 측정한 바 있다.
KR-BERT~\cite{krbert}, KoCharELECTRA~\cite{kocharelectra} 등은 서브워드 vocab 구성 단계에서 Hangul의
구조적 특성을 부분적으로 반영하였으며, Park \emph{et~al.}~\cite{park2020jamo}은 Jamo 단위
word2vec 임베딩이 형태소 단위 임베딩 대비 유사 태스크에서 경쟁력 있는 성능을 보임을 보고하였다.
본 연구는 이러한 계보 위에서 특히 ``Jamo 기저 + BPE merge'' 조합의 토큰화 효율과 학습 영향을
정량적으로 비교한다는 점에 초점을 맞춘다.

\subsection{Tokenizer Evaluation}

토크나이저 평가에서는 전통적으로 \emph{fertility}(원문 단어 대비 토큰 수 비율)와 평균 토큰 길이,
그리고 downstream task 성능이 주로 사용된다~\cite{rust2021multilingualtok}. 최근에는
tokenizer가 LLM의 학습 효율과 다국어 성능에 미치는 영향이 재조명되면서, 언어별 fertility 격차가
다국어 모델의 ``tax''로 작용한다는 점이 여러 차례 지적되었다. 본 연구에서 보고하는
``tiktoken이 한국어에서 비효율적''이라는 결과는 이러한 맥락과 일치한다.

\section{Motivation}
\label{sec:motivation}

한국어에서는 어간과 어미, 형태소의 경계가 종종 음절 \emph{내부}에 위치한다. 두 단어가 결합하면서
발생하는 사잇소리 현상이나, 활용형에서 나타나는 종성 변화가 대표적인 예이다. 문자 단위
토크나이저는 음절을 원자로 취급하므로 이러한 내부 경계를 표현할 수 없다.
그림~\ref{fig:meaning_split}에서와 같이, 의미 단위로 정확히 분절하려면 하나의 문자 블록을
내부에서 쪼개야 하는 상황이 발생한다. 이것이 본 연구에서 Jamo 수준에서 토큰화를 수행하고자 한
핵심 동기이다.

\begin{figure}[h]
\centering
\fbox{\parbox{0.9\columnwidth}{%
\centering
단어: \textbf{노랫말} (가사, lyrics)\\[2pt]
의미 분절: 노래 + ㅅ(사잇소리, interfix) + 말\\[2pt]
문자 단위 분절은 ``노랫''에서 종성 ``ㅅ''을 분리할 수 없음\\[2pt]
Jamo 단위 분절: ㄴㅗㄹㅐ / ㅅ / ㅁㅏㄹ
}}
\caption{문자 내부에서의 의미 경계 예시. Hangul 음절 ``랫''은 ``래''(노래의 일부)와 사잇소리 ``ㅅ''이
결합한 결과이며, Jamo 수준에서만 그 경계를 드러낼 수 있다.}
\label{fig:meaning_split}
\end{figure}

또한 vocab 효율 측면에서도 Jamo 수준 접근은 이점이 있다. 문자 기반 토크나이저는 학습 코퍼스에
등장하는 11,172개 가능한 Hangul 음절을 모두 개별 원자로 취급해야 하지만, Jamo 기반 접근은
67개(19+21+28$-$1, 종성 ``없음'' 제외)의 Jamo만 기저 단위로 두고 나머지 어휘 용량을
빈출 서브워드 병합에 할당할 수 있다. 이는 특히 vocab 크기가 작을 때 한국어 표현력에 유리하게 작용한다.

세 번째로, \emph{희귀 음절(rare syllable)}에 대한 강건성이 개선된다. 문자 기반 토크나이저가 학습 중
관찰하지 못한 음절은 거의 항상 UNK로 처리되거나, 바이트 fallback을 가진 byte-level 토크나이저의
경우 UTF-8 바이트 시퀀스(보통 3바이트)로 파편화된다. 반면 Jamo 기반에서는 해당 음절이
학습 코퍼스에서 관찰되지 않았더라도 그 음절을 구성하는 자모 자체는 코퍼스 전체에서 충분히
빈번하게 등장하므로, 분해된 자모 시퀀스 수준에서 안정적으로 표현된다. 이는 특히 고유명사,
외래어 표기, 고어 등에 대한 분절 품질에 긍정적으로 작용한다.

마지막으로, 공학적 관점에서 Jamo 분해는 \emph{결정론적}이고 \emph{가역적}이다. 외부 형태소 분석기나
사전에 의존하지 않으므로 구현이 단순하고, 배포 환경에 추가 의존성을 가져올 필요가 없다.
모든 현대 Hangul 음절에 대해 분해와 재조합이 $O(1)$ 연산이며, 순방향/역방향 변환이
완전 일대일(bijective)이므로 번역 결과를 다시 한국어로 렌더링할 때 정보 손실이 없다.

\section{Methodology}
\label{sec:method}

\subsection{Data}

\textbf{토크나이저 학습 데이터.} Hugging Face에 공개된 CC-100 한국어 처리
데이터셋(\code{CocoRoF/cc-100-korean-processing})을 사용하였다. 이 데이터셋은 23개의 청크(chunk)로
분할되어 있으며, 대규모 학습 버전(HF 기반)에서는 처음 두 청크(\code{chunk_00}, \code{chunk_01})를
iterator로 순차 주입하여 학습하였다. 소규모 학습 버전(자체 Python 구현)은 공모전 기반 소설 데이터
및 텍스트 파일을 사용하는 경량 구성이다.

\textbf{번역 모델 학습 데이터.} 한국어-영어 병렬 코퍼스는 AI Hub 유형의 ``Tech'' 및 ``Life'' 도메인
\code{ko_en_train_set.json} / \code{ko_en_valid_set.json} 파일에서 추출한
$\langle$ko, en$\rangle$ 쌍을 합쳐 사용하였다.
전처리 단계에서 한국어/영어 중 어느 한 쪽의 길이가 2자 미만인 짝은 제거하였다.

\textbf{전처리.} Jamo 모드(\code{jamo_break=True})에서는 Kss \code{h2hcj}로 한글 음절을
호환 Jamo 시퀀스로 변환한 뒤 BPE 학습/추론에 투입한다. 디코딩 시에는 \code{hcj2h}로
역변환하여 원래의 음절로 복원한다. 문자 모드(\code{jamo_break=False})에서는 별도의 변환 없이
원본 텍스트를 그대로 사용한다.

\textbf{데이터 파이프라인.} Hugging Face \code{datasets.load_dataset}은 \code{datas/huggingface/}
캐시에 요청된 청크를 사전 다운로드한 뒤 로드한다. 각 청크는 평균 수십만\textasciitilde수백만 줄의
웹 크롤링 기반 한국어 텍스트로 구성된다. 본 연구는 전체 23개 청크 중 토크나이저 학습용으로 처음 두
청크(\code{chunk_00}, \code{chunk_01})만 사용하므로 전체 데이터셋을 한 번에 보유할 필요는 없다.
평가용 통계는 학습에 사용되지 않은 \code{chunk_05}에서 5,000개 샘플을 추출하여 측정함으로써
training--evaluation 독립성을 확보하였다.

\subsection{Jamo-based BPE Tokenizer}

본 연구에서는 두 가지 구현의 Jamo 기반 토크나이저를 제공한다.

\textbf{(1) 자체 Python 구현 (\code{JamoBPE}).} 교육적 목적과 알고리즘 이해를 위해 순수 Python으로
BPE를 직접 구현하였다. 알고리즘은 다음의 단계를 따른다:
\begin{enumerate}
    \item 전 코퍼스를 공백 기준으로 토큰화하고, 각 단어를 문자(또는 Jamo) 시퀀스에 \code{</w>}
    (end-of-word) 접미사를 붙인 튜플로 변환한 뒤 빈도를 집계한다 (\code{build_word_freqs}).
    \item 모든 단어에 대해 인접 심볼 쌍의 빈도를 집계한다 (\code{get_pair_stats}).
    \item 최빈 쌍을 병합하여 새로운 심볼을 생성하고 (\code{merge_pair_in_word}),
    코퍼스를 갱신한다 (\code{merge_corpus}).
    \item vocab 크기가 목표치에 도달하거나 최빈 쌍의 빈도가 \code{min_frequency}($=2$) 미만이면 종료한다.
\end{enumerate}
학습 결과는 \code{vocab.json}과 \code{merges.json}으로 직렬화되며, 인코딩 시에는 저장된 merge
순서를 그대로 적용한다. vocab 크기는 기본 32,768로 설정하였다. 단, 소규모 코퍼스 특성상
\code{min_frequency=2}의 종료 조건(\code{jamo_bpe.py} 98행)이 \code{vocab_size} 목표보다 먼저
만족되어, 실제 수렴된 어휘 크기는 약 8{,}904로 형성되었다.

\textbf{(2) Hugging Face 구현 (\code{HFJamoBPE}).} 대규모 데이터로의 확장을 위해 Rust로 구현된
Hugging Face \code{tokenizers} 라이브러리를 래핑하였다. 핵심 설정은 다음과 같다:
\begin{itemize}
    \item \code{model = models.BPE(end_of_word_suffix="</w>")}
    \item \code{pre_tokenizer = pre_tokenizers.Whitespace()}
    \item \code{trainer = trainers.BpeTrainer(vocab_size, min_frequency, special_tokens=[], end_of_word_suffix="</w>")}
    \item BPE 학습 자체에는 특수 토큰을 등록하지 않는다(\code{special_tokens=[]}). 학습 완료 후
    \code{ensure_special_tokens()}로 post-hoc 추가하며, 번역 학습 경로에서는 한국어 측에
    \code{<pad>}, \code{<unk>}만, 영어 GPT-2 측에 \code{<pad>}, \code{<bos>}, \code{<eos>}만
    각각 부여한다. \code{<bos>}/\code{<eos>}는 타겟 생성에만 필요하므로 소스 측에는 불필요하다.
\end{itemize}
본 구현은 vocab 크기 32,768 및 16,384의 두 설정으로 학습하여, vocab 크기 축소가
Jamo/문자 기반 분절 효율 차이에 미치는 영향을 관찰할 수 있도록 하였다.

두 구현 모두 입력 텍스트에 대해 먼저 Hangul 음절을 구성 Jamo로 분해하는 전처리를 수행한 뒤,
해당 Jamo 시퀀스 위에서 BPE 학습 및 토큰화를 진행한다. 디코딩 시에는 토큰을 이어붙이고
\code{</w>}를 공백으로 치환한 뒤 Jamo를 재조합한다.

\textbf{구현 간 차이.} 두 구현은 동일한 알고리즘적 아이디어를 공유하지만 다음의 공학적 차이를 갖는다.
(i) 자체 Python 구현은 단일 스레드 CPython에서 실행되므로 수백만 문장 규모의 학습에는 부적합하다.
따라서 Python 구현은 공모전 기반 소설 텍스트(수 MB 규모)로 학습하고, 개념 증명(proof-of-concept)과
알고리즘 투명성 확보의 역할만 담당한다.
(ii) HF 구현은 Rust로 작성된 native 백엔드와 병렬화된 훈련 파이프라인을 이용하므로, CC-100 수 GB 규모의
데이터에서도 수십 분 이내에 학습이 완료된다. HF 구현은 모든 본 연구의 실제 비교 벤치마크
(표~\ref{tab:stats})와 번역 실험의 주 토크나이저로 사용된다.
(iii) 인코딩 경로의 경우, 자체 구현은 단순한 greedy longest-match 대신 저장된 merge 순서를 반복
적용하는 ``재-병합''(re-merge) 방식을 사용한다. 이는 훈련 시의 분절과 추론 시의 분절이 정확히
일치하도록 보장하지만, 최악의 경우 $O(|\text{merges}| \cdot n)$의 시간 복잡도를 갖는다. HF 구현은
내부적으로 해시 기반의 병합 테이블 조회로 $O(n)$에 가까운 성능을 낸다.

\subsection{Translation Model}

토크나이저가 실제 학습에 미치는 영향을 평가하기 위해, PyTorch 기반으로 표준 Transformer 구조의
한국어-영어 번역기를 구현하였다. 전체 데이터 흐름은 그림~\ref{fig:arch}에 도식화되어 있으며,
구체 아키텍처와 하이퍼파라미터는 표~\ref{tab:model}과 같다.

\begin{figure*}[t]
\centering
\includegraphics[width=0.95\textwidth,keepaspectratio]{figures/transformer.png}
\caption{번역 모델 아키텍처. 한국어 소스(예: ``강아지'')는 embedding 후 4개의 Encoding Layer
(self-attention + feed-forward)를 거쳐 memory 표현을 생성한다. 영어 타겟(예: ``Puppy'')은
embedding 후 4개의 Decoding Layer(causal-masked self-attention + encoder--decoder cross-attention
+ feed-forward)를 거치며, 각 decoding layer는 encoder memory를 cross-attention으로 참조한다.
최종 decoder 출력은 unembedding(선형 사영)을 통해 타겟 어휘 공간으로 투영되고, teacher forcing
하에서 다음 토큰(예: ``Cat?'')과의 cross-entropy loss로 학습된다.}
\label{fig:arch}
\end{figure*}

\begin{table}[h]
\caption{Seq2SeqTransformer 구성.}
\label{tab:model}
\small
\begin{tabular}{ll}
\toprule
구성 요소 & 값 \\
\midrule
$d_\text{model}$ & 512 \\
Attention heads & 8 \\
Encoder layers & 4 \\
Decoder layers & 4 \\
Feed-forward dim. & 1024 \\
Dropout & 0.1 \\
Max sequence length & 128 \\
Positional encoding & sinusoidal \\
Source tokenizer & Jamo-BPE 32k / Full-BPE 32k \\
Target tokenizer & GPT-2 tokenizer (fixed) \\
\midrule
Total parameters & $\approx$ 89M \\
\bottomrule
\end{tabular}
\end{table}

모델은 \code{torch.nn.Transformer}를 기반으로 한다. 입력 임베딩은 $\sqrt{d_\text{model}}$로
스케일링된 뒤 sinusoidal positional encoding과 합산되며, 디코더에서는 상삼각 causal mask로
미래 토큰의 누출을 방지한다. 소스/타겟 모두 패딩 마스크를 적용하여 패딩 토큰이 attention에
영향을 주지 않도록 한다. 최종적으로 디코더 출력을 선형 층(\texttt{generator})을 통해 타겟 어휘 공간으로
사영한 후 teacher forcing 방식으로 학습한다.

\textbf{학습 설정.} 배치 크기 32, 최대 5 epoch, 최적화기는 AdamW($\beta_1{=}0.9$, $\beta_2{=}0.98$,
$\epsilon{=}10^{-9}$), 기준 학습률 $3 \times 10^{-4}$에 inverse-square-root 스케줄(warmup 1,000 step)을
적용하였다. cross-entropy loss에 label smoothing 0.1, gradient clipping $\|g\|_2 \le 1.0$,
CUDA에서는 bfloat16 autocast를 사용한다. 검증 loss가 최저인 시점의 체크포인트를 저장하며,
보고서의 학습 곡선 비교 결과는 실제 실험 기록 기준 4 epoch까지 진행된 run을 사용하였다.

영어 측 타겟 토크나이저는 GPT-2 사전학습 토크나이저(\code{Tokenizer.from_pretrained("gpt2")})를
고정으로 사용하여, 실험 변수를 한국어 측 토크나이저로 제한하였다. 이는 영어 측의 토큰화 품질을
``산업 표준''에 고정시켜, 관찰되는 차이가 한국어 측 토크나이저에서 비롯됨을 분명히 하기 위한
설계이다.

\textbf{데이터 수집 경로 (Data Loading).} 번역 데이터는 \code{build_pairs} 함수가 각 도메인
JSON 파일을 \code{json.load()}로 전체 메모리에 읽은 뒤, $\langle$ko, en$\rangle$ 쌍을 파이썬
리스트에 누적한다(\code{dataset.py} 21, 78행). 한국어/영어 중 길이 2자 미만인 짝은 이 단계에서
필터링된다. 각 배치는 \code{make_collate_fn}이 생성한 collate 함수에 의해 소스/타겟 시퀀스가
각각 \code{<pad>}로 right-pad되어 배치 텐서로 변환되고, padding 마스크와 teacher-forcing용
\code{tgt_in}/\code{tgt_out} (right/left shifted) 텐서가 함께 생성된다. \code{DataLoader}는
훈련 시 \code{shuffle=True}와 \code{drop_last=True}로, 검증 시 \code{shuffle=False}로 구성된다.

\textbf{로그 및 체크포인트.} 매 학습 run은 \code{runs/{kr_tokenizer}/{timestamp}/} 경로에
\code{log.jsonl}과 \code{best.pt}를 생성한다. 100 스텝마다 step-level loss, learning rate,
유효 토큰 수가 JSONL 형식으로 기록되며, 매 epoch 종료 시 train/valid loss와 5개 검증 샘플에 대한
greedy 번역 결과도 함께 기록된다. 이 로그는 \code{src/scripts/plot_loss.py}를 통해 매끄러운 EMA
스무딩 곡선으로 시각화된다.

\subsection{Tokenizer Benchmarking}

토큰화 효율을 측정하기 위해 \code{src/scripts/tokenizer_stats.py}를 통해 CC-100 한국어
\code{chunk_05}의 최대 5,000개 문장에 대해 각 토크나이저로 인코딩한 뒤 토큰 개수 분포를
집계하였다. 측정 지표는 샘플 수, 총 토큰 수, 평균(mean), 중앙값(median), 표준편차(stdev),
분산(var), 최솟값(min), 95-분위수(p95), 최댓값(max)이다. 비교 대상으로는 자체 Python 구현
Jamo/Full 32k, HF 구현 Jamo/Full 32k 및 16k, 그리고 OpenAI tiktoken \texttt{o200k\_base}(GPT-4o 인코딩)를
포함하였다.

\subsection{Demo}

연구 결과물의 접근성을 높이기 위해 Streamlit 기반의 데모를 제작하였다
(\code{src/main.py}). 사용자는 (i) Dataset 규모 (\code{small} / \code{large} / \code{large_16k}),
(ii) Mode (\code{jamo} / \code{full})의 두 축을 선택한 뒤 한국어 문장을 입력하면
분절 결과를 실시간으로 확인할 수 있고, \code{large} 설정에서는 해당 토크나이저로 학습된
번역 체크포인트를 로드하여 greedy decoding으로 영어 번역까지 함께 출력한다.

데모 앱은 \code{@st.cache_resource}를 사용해 토크나이저와 모델 체크포인트를 세션에 한 번만 로드하므로
사용자 상호작용의 지연을 최소화하였다.

\subsection{Implementation Details}

전체 프로젝트는 Python 3.12, PyTorch ($\ge$2.x), CUDA 12.8 환경에서 수행되었다.
주요 의존성은 \code{tokenizers} $\ge$0.22, \code{kss} $\ge$6.0, \code{datasets} $\ge$4.8,
\code{tiktoken} $\ge$0.12, \code{streamlit} $\ge$1.56이다. 모든 실험은 단일 CUDA GPU에서
수행되었으며, bfloat16 autocast를 통해 메모리와 처리량을 최적화하였다.

파일 구성은 다음과 같다:
\begin{itemize}
    \item \code{src/tokenizer/jamo_bpe.py}, \code{src/tokenizer/hf_jamo_bpe.py} — 토크나이저 구현.
    \item \code{src/translation/model.py} — \code{Seq2SeqTransformer} 정의 및 greedy decode.
    \item \code{src/translation/train.py} — 학습 루프, 로그, 체크포인트 저장.
    \item \code{src/translation/dataset.py} — JSON 기반 병렬 데이터 로딩과 collate.
    \item \code{src/scripts/train_hf_tokenizer.py} — HF 토크나이저 대규모 학습 스크립트.
    \item \code{src/scripts/tokenizer_stats.py} — 토크나이저 벤치마크 스크립트.
    \item \code{src/scripts/plot_loss.py} — 학습 loss 곡선 시각화.
    \item \code{src/main.py} — Streamlit 데모 엔트리포인트.
\end{itemize}

\section{Experimental Setup}
\label{sec:setup}

토크나이저 비교 실험은 표~\ref{tab:setup}과 같이 6개의 한국어 토크나이저 변형에 대해 수행하였다.
번역 학습 실험은 이 중 \texttt{hf\_jamo}와 \texttt{hf\_full}(32k, CC-100 2 청크 학습) 두 조건에서만
수행하였다.

\begin{table}[h]
\caption{실험에 사용된 토크나이저 변형.}
\label{tab:setup}
\small
\begin{tabular}{llll}
\toprule
Key & 구현 & Mode & Vocab \\
\midrule
\texttt{local/jamo}      & Python BPE & Jamo & 32k (실제 8{,}904) \\
\texttt{local/full}      & Python BPE & Char & 32k (실제 8{,}904) \\
\texttt{hf/jamo}         & HF BPE     & Jamo & 32k \\
\texttt{hf/full}         & HF BPE     & Char & 32k \\
\texttt{hf/jamo\_16k}    & HF BPE     & Jamo & 16k \\
\texttt{hf/full\_16k}    & HF BPE     & Char & 16k \\
\texttt{tiktoken}        & OpenAI     & Byte-BPE & \texttt{o200k\_base} \\
\bottomrule
\end{tabular}
\end{table}

\section{Results}
\label{sec:results}

\subsection{Tokenization Efficiency}

CC-100 한국어 \texttt{chunk\_05}에서 추출한 5,000개 문장에 대한 토큰 길이 분포 통계를
표~\ref{tab:stats}에 정리하였다.

\begin{table}[h]
\caption{토크나이저별 토큰 수 통계 (CC-100 ko / \texttt{chunk\_05}, 5{,}000 samples).}
\label{tab:stats}
\small
\begin{tabular}{lrrrr}
\toprule
Tokenizer & Total & Mean & Min & p95 \\
\midrule
\texttt{local/jamo} 32k    & 197{,}285 & 39.46 & 4  & 125 \\
\texttt{local/full} 32k    & 292{,}736 & 58.55 & 10 & 197 \\
\texttt{hf/jamo} 32k       & 124{,}532 & 24.91 & 2  & 83  \\
\texttt{hf/full} 32k       & 126{,}119 & 25.22 & 2  & 83  \\
\texttt{hf/jamo\_16k}      & 137{,}542 & 27.51 & 3  & 91  \\
\texttt{hf/full\_16k}      & 145{,}554 & 29.11 & 3  & 95  \\
\texttt{tiktoken/o200k}    & 182{,}342 & 36.47 & 2  & 120 \\
\bottomrule
\end{tabular}
\end{table}

통계 결과에서 두 가지 경향이 명확하게 관찰된다.
\begin{itemize}
    \item \textbf{소설 코퍼스로 학습한 토크나이저에서 Jamo 우위가 극대화된다.} 공모전 기반
    소설 텍스트(수 MB 규모)로 학습된 자체 Python 구현 기준, \texttt{local/jamo}의 평균 토큰 수는
    39.46으로 \texttt{local/full}의 58.55 대비 약 32.6\% 적다. 평가는 CC-100 \code{chunk_05}에서
    이루어지므로, 학습-평가 도메인 불일치 상황에서 문자 기반 토크나이저는 미관찰 음절을 병합할
    기회가 부족한 반면, Jamo 기반은 초성/중성/종성 수준의 풍부한 공유 통계로 안정적인 분절을
    유지하기 때문으로 해석된다.
    \item \textbf{Vocab 크기가 작을수록 Jamo의 상대적 이점이 커진다.} HF 구현 기준, vocab 32k에서는
    Jamo(24.91)와 Full(25.22)의 평균 토큰 수 차이가 약 1.3\%에 불과하지만,
    16k로 축소하면 Jamo(27.51)는 Full(29.11) 대비 약 5.5\% 적다. vocab 용량이 제한될 때
    Jamo 기반이 더 효율적으로 한국어를 표현함을 시사한다.
    \item \textbf{다국어 토크나이저는 한국어에서 비효율적이다.} tiktoken \texttt{o200k\_base}의 평균 토큰 수는
    36.47로, 본 연구의 한국어 전용 32k 토크나이저(Jamo, Full 모두) 대비 약 45\% 이상 더 많은 토큰을
    사용하였다. 이는 다국어 모델이 한국어에 제한된 어휘 예산만을 할당하는 구조적 결과이다.
\end{itemize}

흥미로운 점은 \emph{최소 토큰 수(min)} 지표에서 HF 32k 토크나이저와 tiktoken이 모두 2로 동일하다는
것이다. 즉 매우 짧은 입력에서는 tiktoken의 긴 다국어 서브워드가 유리하게 작용하는 경우가
존재한다. 반면 본 연구의 16k 토크나이저는 최솟값이 3으로, 충분히 짧은 입력을 단일 또는 2-토큰으로
커버하지 못한다는 trade-off를 보였다. 이 격차는 vocab 예산이 줄어들수록 ``가장 짧은 표현'' 영역에
할당되는 고빈도 서브워드의 수가 감소하기 때문으로 해석된다.

\subsection{Distribution Analysis}

평균 이외의 분포 특성 또한 토크나이저 선택에 중요한 함의를 제공한다.
표~\ref{tab:stats_full}은 표준편차와 95-분위수를 포함한 확장된 통계이다.

\begin{table}[h]
\caption{토크나이저별 토큰 수의 분포 특성 (CC-100 ko / \texttt{chunk\_05}).}
\label{tab:stats_full}
\small
\begin{tabular}{lrrrr}
\toprule
Tokenizer & Median & Stdev & p95 & Max \\
\midrule
\texttt{local/jamo} 32k    & 24 & 45.31 & 125 & 627 \\
\texttt{local/full} 32k    & 32 & 73.11 & 197 & 1293 \\
\texttt{hf/jamo} 32k       & 15 & 30.04 & 83  & 523 \\
\texttt{hf/full} 32k       & 15 & 30.44 & 83  & 526 \\
\texttt{hf/jamo\_16k}      & 16 & 33.21 & 91  & 570 \\
\texttt{hf/full\_16k}      & 17 & 35.13 & 95  & 609 \\
\texttt{tiktoken/o200k}    & 21 & 44.67 & 120 & 781 \\
\bottomrule
\end{tabular}
\end{table}

주목할 부분은 다음과 같다.
\begin{itemize}
    \item \textbf{꼬리(tail) 특성.} \texttt{local/full}의 최댓값은 1,293 토큰으로, 같은 데이터를
    \texttt{hf/full} 32k는 최대 526 토큰으로 처리한다. 이는 작은 학습 코퍼스에서 훈련된 문자 기반
    토크나이저가 긴 문장 처리 시 극단적으로 비효율적일 수 있음을 시사하며, 실제 서비스 환경에서는
    context window 초과 위험과 직결된다.
    \item \textbf{표준편차 축소.} HF 32k 체계에서 Jamo(30.04)와 Full(30.44)의 stdev가 거의 동일한
    반면, 16k로 축소하면 Jamo(33.21)가 Full(35.13) 대비 약 5.5\% 낮은 변동성을 보인다. 평균뿐 아니라
    분산 측면에서도 Jamo 기반이 더 일관된 분절을 제공함을 뜻한다.
    \item \textbf{중앙값 기반 비교.} 평균은 긴 꼬리에 민감하므로, 중앙값(median) 기준으로도
    비교하였다. \texttt{hf/jamo\_16k}(16) vs \texttt{hf/full\_16k}(17)에서 Jamo 기반의 우위가
    꼬리 효과가 아닌 전반적 분포의 중앙 이동에서 비롯됨을 확인할 수 있다.
\end{itemize}

\subsection{Training Loss}

동일 구조(표~\ref{tab:model})의 Transformer 번역기를 \texttt{hf\_full} 32k와 \texttt{hf\_jamo} 32k
토크나이저로 각각 학습시킨 결과, 학습 단계별 loss 곡선은 그림~\ref{fig:loss}와 같이 서로
매우 유사하게 나타났다. 지수 이동평균($\alpha=0.9$)으로 스무딩한 곡선에서도 두 run의 차이는 일관되게
작았으며, 최종 validation loss의 상대 차이는 약 1\% 수준으로 본 실험 규모에서는 오차 범위 내로
해석할 수 있는 차이이다.

\begin{figure}[h]
\centering
\includegraphics[width=0.95\columnwidth,keepaspectratio]{src/scripts/plots/combined_smooth.png}
\caption{\texttt{hf\_full} vs \texttt{hf\_jamo} 학습 loss 비교 (EMA 0.9 스무딩).}
\label{fig:loss}
\end{figure}

이는 토큰화 \emph{효율}의 차이가 loss 관점에서는 거의 드러나지 않음을 뜻한다. 두 토크나이저가
비슷한 수준의 시퀀스 길이를 생성하고 있고(표~\ref{tab:stats} 참조), Transformer가 결국
어떤 분절 방식이든 동등하게 표현력을 학습해 내기 때문인 것으로 추정된다.

추가로, loss 수치는 토큰 단위로 계산되므로 단순한 loss 비교는 토크나이저 간 ``동등'' 비교가 아님에
주의해야 한다. Jamo 기반은 문장당 토큰 수가 약간 적어 한 토큰당 정보량이 더 크고, 따라서 절대 loss가
동일하더라도 실제 ``문장당 perplexity'' 혹은 번역 품질 관점에서는 차이가 존재할 수 있다. 본 실험에서는
시퀀스 길이 차이가 작아 이러한 보정의 영향이 제한적이지만, 더 큰 vocab 차이가 있는 조건에서는
bits-per-byte(bpb) 혹은 character-level loss 등 토크나이저 독립적 지표로의 환산이 필요할 것이다.

\subsection{Qualitative Analysis: Meaning-aligned Segmentation}

정량 결과 외에, 의미 경계와 토큰 경계의 일치 측면에서 흥미로운 정성적 결과를 확인하였다.
단어 ``노랫말''(``lyrics'')에 대해, 문자 기반 토크나이저는 이를 세 개의 토큰으로 분절하였으나
앞의 두 토큰은 그 자체로 독립적인 의미를 갖지 않았고 오직 마지막 토큰만이 ``말''(words)에
해당하는 의미를 가졌다.

반면 Jamo 기반 토크나이저는 동일하게 세 개의 토큰으로 분절하였지만, 각 토큰이 각각
``노래''(music/song), 사잇소리 ``ㅅ''(두 단어를 연결하는 interfix), 그리고 ``말''(words)에
대응되도록 분리되었다. 즉, 형태소 경계와 일치하는 의미 중심의 분절이 이루어진 것이다.
이러한 정렬은 모델이 의미 단위에 대해 더 일관된 표현(representation)을 학습할 수 있는 조건을
제공한다고 볼 수 있다.

이 외에도 활용 어미 경계의 예시를 살펴보자. 동사 ``먹다''의 존댓말 활용형 ``먹었어요''는
``먹(어간) + 었(과거 선어말 어미) + 어요(종결 어미)''로 분석되는데, 여기서 ``었''의 종성 ``ㅆ''는
본래 어간의 일부가 아니라 과거 시제 형태소의 일부이다. 문자 기반 토크나이저는 ``었''을 하나의
원자로 고정하므로 이러한 형태소 경계를 드러낼 기회가 원리적으로 차단되지만, Jamo 기반은
ㅆ를 별도의 자모로 관찰하므로 ``ㅆ + 어요'' 혹은 ``ㅆ어요''와 같은 형태소 기반 분절을 자연스럽게
학습할 여지를 갖는다. 이는 한국어와 같이 교착어적 성격을 지닌 언어에서 특히 유효하다.

\subsection{Real-world Examples}

실제 한국어 가사 전문(全文)에 대한 적용 사례로, BTS의 ``봄날''(Spring Day)과 IU의
``밤편지'' 두 곡을 각각 네 토크나이저(Small/Large $\times$ Jamo/Character)에 투입하여
총 토큰 수를 측정하였다. 결과는 표~\ref{tab:lyrics}와 같다.

\begin{table}[h]
\caption{실 가사 전문(全文) 토큰화 결과. \textbf{Small}은 소설 코퍼스(공모전 기반, 수 MB 규모)로
학습한 자체 Python BPE(\code{local/*}, 수렴 vocab $\approx$ 8{,}904), \textbf{Large}는 CC-100
한국어 2 청크로 학습한 HF BPE 32k(\code{hf/*})를 가리킨다.}
\label{tab:lyrics}
\small
\begin{tabular}{llrr}
\toprule
곡 & 데이터 & Jamo Based & Character Based \\
\midrule
BTS -- 봄날       & Small & 530 & 953 \\
                  & Large & 436 & 443 \\
\midrule
IU -- 밤편지      & Small & 165 & 352 \\
                  & Large & 145 & 148 \\
\bottomrule
\end{tabular}
\end{table}

두 곡 모두에서 CC-100 벤치마크(표~\ref{tab:stats})의 패턴이 그대로 재현된다.
\begin{itemize}
    \item \textbf{Small 조건}(소설 코퍼스 학습 Python BPE)에서 격차가 크다. 봄날에서 Jamo
    530 vs Char 953(약 44.4\% 절감), 밤편지에서 Jamo 165 vs Char 352(약 53.1\% 절감)으로,
    Jamo 기반이 훨씬 짧은 시퀀스를 생성함이 일관되게 관찰된다. 이는 도메인 한정 소설 코퍼스가
    일반 한국어(가사)에서 등장하는 음절을 충분히 포괄하지 못해 문자 기반 쪽이 음절 단위 병합을
    학습할 기회가 적었던 결과로 해석된다.
    \item \textbf{Large 조건}(CC-100 2 청크 학습 HF BPE 32k)에서는 격차가 크게 축소된다.
    봄날에서 436 vs 443(약 1.6\% 차이), 밤편지에서 145 vs 148(약 2.0\% 차이)으로 좁혀진다.
    웹 크롤링 기반 CC-100은 가사와 어휘/표현의 겹침이 크고, 수십만 줄 규모로 모든 흔한 Hangul
    음절을 충분히 관찰하므로 문자 기반에도 풍부한 병합 기회가 주어지기 때문이다.
    \item \textbf{곡 간 일관성}. 두 곡 모두 Jamo 기반이 Character 기반보다 짧거나 동등한
    토큰 시퀀스를 생성하며, 열세인 사례는 관찰되지 않았다.
\end{itemize}

밤편지의 Small 조건 절감률(53.1\%)이 봄날(44.4\%)보다 크게 나타나는 것은 가사 내 반복
후렴구 비중 및 고유 음절 다양성 차이에 기인한 것으로 보이며, 곡 단위 정성적 관찰에 해당한다.
두 사례 모두 \S\ref{sec:results}의 CC-100 평균 토큰 수 경향(``Small에서 Jamo 우위 극대화, Large에서 근사'')과
정합한다.

\section{Discussion}
\label{sec:discussion}

실험 결과는 다음과 같이 종합할 수 있다.
\begin{itemize}
    \item \textbf{토큰화 효율성 (수치적).} Jamo 기반 토크나이저는 동일하거나 더 작은 vocab 크기에서도
    평균적으로 더 적은 토큰 수로 한국어 문장을 표현할 수 있었다. 이 이점은 코퍼스 규모가 작거나
    vocab 예산이 제한될수록 확대된다.
    \item \textbf{의미 일치성 (구조적).} 사잇소리와 같이 음절 내부에 형태소 경계가 존재하는 한국어 고유의
    언어 현상이 분절 경계에 자연스럽게 반영되는 사례가 확인되었다. 이는 문자 기반 BPE에서는
    원리적으로 불가능한 분절이다.
    \item \textbf{학습 성능 (다운스트림).} 본 연구의 실험 규모(약 89M 파라미터 Transformer;
    \code{src_vocab=32,770}, \code{tgt_vocab=50,260}, \code{d_model=512}, 4+4 layer이며
    embedding과 output projection이 전체 파라미터의 약 75\%를 차지,
    2 청크 분량 학습 데이터, 최대 5 epoch)에서는 두 토크나이저 간 학습 곡선의 차이가 미미하였다.
    따라서 Jamo 기반 분절의 구조적 이점이 실제 번역 품질(BLEU, chrF, COMET 등)에 유의미하게
    반영되는지는 본 규모에서는 단정할 수 없으며, 향후 더 큰 모델과 더 긴 학습 스케줄에서의
    검증이 필요하다.
    \item \textbf{Trade-off.} Jamo 분해는 시퀀스를 \emph{음절 단위로 보았을 때} 길게 만든다
    (1 음절 = 2\textendash3 Jamo). 병합을 통해 이 길이를 회복해야 하므로, 단순한 ``jamo = 언제나 더 짧은
    시퀀스''라는 결론은 성립하지 않는다. 실제로 HF 32k에서 Jamo와 Full의 평균 토큰 수는 거의
    동일하며, 우위는 vocab/데이터 예산이 제한될 때 주로 드러난다.
    \item \textbf{실무적 함의.} 다국어 LLM 서비스에서 한국어 fertility가 과도하게 높은 문제를 완화하기
    위한 한 가지 실용적 접근은, 다국어 공통 vocab에 Jamo 기반 merge를 포함하는 것이다.
    본 연구의 결과는 Jamo 기반이 문자 기반보다 평균 토큰 수에서 결코 열세가 아니라는 점을 보이므로
    (32k에서 동등, 16k에서 우위), 다국어 vocab에 Hangul 음절을 전부 포함시키는 대신 자모를 기저로 두는
    선택이 비용-효율 측면에서 유리할 수 있다.
    \item \textbf{학습 안정성.} 학습 초기 단계(warmup 구간)에서 두 run 모두 loss 곡선이 유사한 궤적을
    그렸고, 학습 종료 시점의 validation loss도 1\% 이내에서 수렴하였다. 이는 Jamo 분해가 모델
    수렴성(convergence stability)에 부정적 영향을 주지 않음을 의미하며, 기존 파이프라인을
    교체하는 비용이 작다는 실무적 장점을 시사한다.
\end{itemize}

\section{Limitations}
\label{sec:limitations}

본 연구는 몇 가지 한계를 지닌다.
\begin{itemize}
    \item \textbf{실험 규모.} 번역 모델은 약 89M 파라미터 규모의 Transformer이며
    (embedding/output projection이 대부분을 차지), 학습 데이터는
    CC-100 2 청크 + AI-Hub 일부 도메인에 제한되었다. 이 규모에서 관찰된 loss 동등성은 더 큰 모델이나
    긴 학습에서 다른 양상을 보일 수 있다.
    \item \textbf{단일 태스크.} 한영 번역이라는 단일 downstream 태스크에서만 평가하였으며,
    언어 모델링, 요약, 분류 등 다른 태스크로의 일반화는 검증되지 않았다.
    \item \textbf{평가 지표.} BLEU, chrF++, COMET 등 다운스트림 품질 지표는 본 보고서에 포함되지
    않았으며, 학습 loss 수렴성만을 비교하였다. 따라서 ``토큰화 효율 개선이 번역 품질로 이어지는가''라는
    질문에는 직접적인 답을 제시하지 못한다.
    \item \textbf{한국어 표기 범위.} Jamo 분해는 현대 Hangul 음절에 대해 완전하지만, 고어 Hangul
    (\texttt{U+1100--U+11FF}의 전체 범위 활용)이나 중국어/일본어가 섞인 혼용 표기 등 경계
    케이스에서는 Kss의 동작에 의존하며, 일부 특수 케이스는 본 연구에서 검증되지 않았다.
    \item \textbf{영어 측 토크나이저 고정.} 타겟 측은 GPT-2 토크나이저로 고정되어 있어, 타겟 토큰화가
    번역 품질에 미치는 영향은 본 연구 범위 밖이다.
    \item \textbf{Greedy decoding.} 번역 예시는 greedy decoding으로 생성되었으며, beam search나
    샘플링 기반 디코딩에서의 차이는 측정하지 않았다.
\end{itemize}

\section{Future Work}
\label{sec:future}

향후 다음의 방향으로 연구를 확장할 수 있다.
\begin{itemize}
    \item \textbf{더 큰 규모의 검증.} 본 연구의 89M 규모 Transformer 대신 수백 M 이상 규모의 모델과
    전체 CC-100 한국어 코퍼스(23 청크) 및 AI-Hub 전 도메인을 활용한 번역 실험을 통해,
    Jamo 기반 토크나이저의 이점이 BLEU, chrF++~\cite{popovic2017chrf}, COMET~\cite{rei2020comet} 등
    다운스트림 지표에 반영되는지 검증한다.
    \item \textbf{다양한 태스크로 확장.} 번역뿐 아니라 요약, 질의응답, 언어 모델링, 형태소 분석 등
    다양한 태스크에서 평가한다. 특히 형태소 분석과 같이 경계 민감 태스크에서는 Jamo 기반 분절의
    이점이 상대적으로 더 크게 드러날 것으로 기대된다.
    \item \textbf{Cross-lingual 일반화.} Jamo 기반 접근을 일본어의 가나/한자 혼용이나 중국어의 부수(radical)
    수준 분해, 베트남어의 결합형 모음/성조 기호 등 다른 동아시아 언어에 유사하게 적용할 수 있는지
    탐구한다.
    \item \textbf{하이브리드 토크나이저.} 형태소 분석기와의 결합, 또는 Jamo-level BPE와
    character-level BPE를 혼합하여, 빈출 음절은 음절 단위로 유지하되 드물고 복합적인 음절만 Jamo 단위로
    분해하는 하이브리드 토크나이저를 설계한다. 이는 평균 시퀀스 길이를 더 줄이면서도 의미 경계 표현력을
    보존할 수 있을 것으로 기대된다.
    \item \textbf{토큰 분포 및 vocab 진단.} BPE merge 순서, Jamo/음절 단위 unigram 확률,
    per-morpheme coverage 등 vocab 수준 진단 지표를 추가하여, 왜 특정 조건에서 Jamo 기반이
    더 유리한지에 대한 정량적 설명을 제공한다.
    \item \textbf{다양한 디코딩 전략.} 본 연구는 greedy decoding만을 사용하였으나, beam search,
    top-$k$/top-$p$ sampling, length-penalty 등 다양한 디코딩 전략 하에서 Jamo 기반 토크나이저가
    생성하는 번역의 품질과 유창성(fluency)을 비교한다.
    \item \textbf{도메인 전이 및 robustness.} 학습 도메인(뉴스/웹)과 다른 도메인(구어체, SNS, 의료 등)에
    대한 토큰화 효율을 측정하여, Jamo 기반 접근이 도메인 외 일반화(out-of-distribution generalization)에
    어떤 영향을 주는지 조사한다.
    \item \textbf{Jamo-level sharing.} 다국어 토크나이저에 Jamo를 기저로 포함시키고, 중국어 부수,
    일본어 가나 등과 어휘를 공유하도록 설계함으로써 여러 동아시아 언어에 걸쳐 vocab을 공유하는 접근을
    탐구한다.
    \item \textbf{Attention/표현 공간 분석.} 동일한 Transformer가 Jamo 기반 vs 문자 기반 입력을 받았을 때
    attention pattern, 임베딩 유사도, 내부 표현의 클러스터링이 어떻게 달라지는지 정성/정량적으로
    비교한다. 이는 ``의미 일치성''이라는 주장을 내부 표현 수준에서 검증하는 데 기여할 수 있다.
\end{itemize}

\section{Conclusion}
\label{sec:conclusion}

본 연구에서는 Hangul의 구조적 특성을 반영하여, 문자 단위가 아닌 Jamo 단위에서 BPE를 적용하는
한국어 토크나이저를 제안하고 구현하였다. 자체 Python 구현(\texttt{JamoBPE})과
Hugging Face 기반 구현(\texttt{HFJamoBPE}) 두 버전을 통해 vocab 크기 32k 및 16k의 다양한 설정에서
실험을 수행하였으며, 동일 구조의 Transformer 기반 한영 번역기를 활용해 실제 학습 곡선을
비교하였다. 그 결과, Jamo 기반 토크나이저는 (i) 동일 vocab에서 평균적으로 더 적은 토큰 수로
한국어를 표현하며, (ii) 작은 vocab 또는 작은 코퍼스에서 그 이점이 확대되고,
(iii) 음절 내부에 걸친 형태소 경계를 표현할 수 있다는 구조적 장점을 제공함을 확인하였다.
다만 본 규모의 번역 실험에서는 학습 loss 측면의 차이는 미미하였으며, 이 구조적 이점이
다운스트림 성능으로 이어지는지는 더 큰 규모에서의 추가 검증이 필요하다. 본 접근은 앞으로 보다 큰
모델과 다양한 태스크에서 검증될 때 그 유용성이 더욱 분명해질 것으로 기대된다.

%% 참고문헌
\begin{thebibliography}{9}

\bibitem{sennrich2016bpe}
R.~Sennrich, B.~Haddow, and A.~Birch,
``Neural machine translation of rare words with subword units,''
in \emph{Proc. of ACL}, 2016.

\bibitem{vaswani2017attention}
A.~Vaswani, N.~Shazeer, N.~Parmar, J.~Uszkoreit, L.~Jones, A.~N.~Gomez, Ł.~Kaiser, and I.~Polosukhin,
``Attention is all you need,''
in \emph{Proc. of NeurIPS}, 2017.

\bibitem{huggingface_tokenizers}
Hugging Face,
``Tokenizers: Fast state-of-the-art tokenizers optimized for research and production,''
\url{https://github.com/huggingface/tokenizers}.

\bibitem{kss}
H.~Park,
``Kss: A Toolkit for Korean sentence segmentation,''
\url{https://github.com/hyunwoongko/kss}.

\bibitem{tiktoken}
OpenAI,
``tiktoken: A fast BPE tokeniser for use with OpenAI's models,''
\url{https://github.com/openai/tiktoken}.

\bibitem{wu2016google}
Y.~Wu \emph{et al.},
``Google's neural machine translation system: Bridging the gap between human and machine translation,''
arXiv:1609.08144, 2016.

\bibitem{kudo2018sentencepiece}
T.~Kudo and J.~Richardson,
``SentencePiece: A simple and language independent subword tokenizer and detokenizer for neural text processing,''
in \emph{Proc. of EMNLP (System Demonstrations)}, 2018.

\bibitem{conneau2020cc100}
A.~Conneau \emph{et al.},
``Unsupervised cross-lingual representation learning at scale,''
in \emph{Proc. of ACL}, 2020.

\bibitem{popovic2017chrf}
M.~Popović,
``chrF++: words helping character n-grams,''
in \emph{Proc. of WMT}, 2017.

\bibitem{rei2020comet}
R.~Rei, C.~Stewart, A.~C.~Farinha, and A.~Lavie,
``COMET: A neural framework for MT evaluation,''
in \emph{Proc. of EMNLP}, 2020.

\bibitem{kudo2018sp}
T.~Kudo,
``Subword regularization: Improving neural network translation models with multiple subword candidates,''
in \emph{Proc. of ACL}, 2018.

\bibitem{provilkov2020bpedropout}
I.~Provilkov, D.~Emelianenko, and E.~Voita,
``BPE-Dropout: Simple and effective subword regularization,''
in \emph{Proc. of ACL}, 2020.

\bibitem{radford2019gpt2}
A.~Radford, J.~Wu, R.~Child, D.~Luan, D.~Amodei, and I.~Sutskever,
``Language models are unsupervised multitask learners,''
OpenAI Technical Report, 2019.

\bibitem{mecabko}
T.~Kudo \emph{et al.},
``MeCab: Yet Another Part-of-Speech and Morphological Analyzer,''
\url{https://taku910.github.io/mecab/}; MeCab-ko fork for Korean.

\bibitem{krbert}
S.~Lee \emph{et al.},
``KR-BERT: A small-scale Korean-specific language model,''
arXiv:2008.03979, 2020.

\bibitem{kocharelectra}
J.~Park,
``KoCharELECTRA: Character-level Korean ELECTRA,''
\url{https://github.com/monologg/KoCharELECTRA}.

\bibitem{park2020jamo}
S.~Park \emph{et al.},
``Sub-character level embeddings for Korean,''
in \emph{Proc. of LREC / 관련 학술대회}, 2020.

\bibitem{rust2021multilingualtok}
P.~Rust, J.~Pfeiffer, I.~Vulić, S.~Ruder, and I.~Gurevych,
``How good is your tokenizer? On the monolingual performance of multilingual language models,''
in \emph{Proc. of ACL}, 2021.

\end{thebibliography}

\end{document}
